\colorlet{appendixHeader}{CadetBlue!20}
\colorlet{appendixStripe}{gray!10}
\colorlet{appendixInk}{black}
\colorlet{appendixRule}{black}
\newcommand{\appendixHeading}[1]{\textcolor{appendixInk}{\textbf{#1}}}

\section{Overview of the Training Algorithm}
\label{app:algorithm}

Algorithm~\ref{alg:collaboration} summarizes the two training stages: protocol distillation and pyramid supervision. Let $\mathcal{M}=\{M_1,\ldots,M_N\}$ denote the heterogeneous agent pool. The teacher model $M^\star$, selected as the largest model in this pool by parameter count, generates trajectories for protocol distillation; we use Qwen3-8B in our experiments. The original large model $M_L$ (Qwen3-14B) serves as the inspector $\mathcal{J}$ during training on all five datasets but is not deployed during collaborative inference.

Let $\mathcal{Q}_{\mathrm{distill}}$ and $\mathcal{Q}_{\mathrm{pyramid}}$ denote the task inputs used for protocol distillation and pyramid supervision, respectively. For each task $q$, we use the trajectory notation $\tau=(z_1,\ldots,z_T)$ with $z_t=(i_t,x_t,y_t,a_t)$ defined in Section~\ref{sec:collaborative_reasoning}, where $x_t$ contains $q$ and the preceding interaction history. The training target $o_t$ is the full serialized response at step $t$, including its reasoning, answer, and interaction fields. The distribution $P_{\mathcal{M}}(\tau\mid q)$ is induced by the agent pool and its interaction protocol. During protocol distillation, a permutation $\pi$ is sampled uniformly from $S_N$, the set of permutations of the $N$ agent identities, and used to remap identities consistently throughout each teacher trajectory. The retained, remapped steps form the supervised dataset $\mathcal{D}_{\mathrm{distill}}$. During pyramid supervision, the inspector assigns a score $s_t^{\mathrm{ins}}$ to each step. Task feedback and inspection scores determine the training weight $r_t$ according to Appendix~\ref{app:training}. For each agent $M_i$, the dataset $\mathcal{D}_i$ contains $(x_t,o_t,r_t)$ from steps with $i_t=i$, so each model is optimized using only its own actions.

\begingroup
\makeatletter
\let\appendixOriginalRuled\fs@ruled
\renewcommand{\fs@ruled}{%
  \appendixOriginalRuled
  \def\@fs@capt##1##2{%
    \begingroup
    \setlength{\fboxsep}{2pt}%
    \noindent\colorbox{appendixHeader}{%
      \parbox{\dimexpr\linewidth-2\fboxsep\relax}{%
        \color{appendixInk}\floatc@ruled{##1}{##2}}}\par
    \endgroup
  }%
  \def\@fs@pre{{\color{appendixRule}\hrule height.8pt depth0pt}\kern2pt}%
  \def\@fs@mid{\kern2pt{\color{appendixRule}\hrule}\kern2pt}%
  \def\@fs@post{\kern2pt{\color{appendixRule}\hrule}}%
}
\makeatother
\begin{algorithm}[H]
\caption{Protocol Distillation and Pyramid Supervision for Heterogeneous Collaboration}
\label{alg:collaboration}
\small
\begin{algorithmic}[1]
\Require Agent pool $\mathcal{M}$; largest model in the pool $M^\star$; task collections $\mathcal{Q}_{\mathrm{distill}},\mathcal{Q}_{\mathrm{pyramid}}$; inspector $\mathcal{J}$; task-specific reward configuration
\State Construct homogeneous teacher MAS $\mathcal{M}_{\mathrm{teacher}}=\{M_1^\star,\ldots,M_N^\star\}$
\State \appendixHeading{Stage I: Protocol Distillation}
\State $\mathcal{D}_{\mathrm{distill}}\gets\emptyset$
\For{$q\in\mathcal{Q}_{\mathrm{distill}}$}
  \State Roll out teacher trajectory $\tau\sim P_{\mathcal{M}_{\mathrm{teacher}}}(\tau\mid q)$
  \State Retain examples satisfying the task-specific data filters
  \State Sample identity permutation $\pi\sim\mathrm{Uniform}(S_N)$
  \State Apply $\pi$ consistently to all agent identities in $\tau$
  \State Add remapped steps to $\mathcal{D}_{\mathrm{distill}}$
\EndFor
\For{$i=1,\ldots,N$}
  \State Distill protocol behavior into $M_i$ using its corresponding steps in $\mathcal{D}_{\mathrm{distill}}$
\EndFor
\State \appendixHeading{Stage II: Pyramid Supervision}
\For{each training iteration}
  \State $\mathcal{D}_1,\ldots,\mathcal{D}_N\gets\emptyset$
  \For{$q\in\mathcal{Q}_{\mathrm{pyramid}}$}
    \State Collect or retrieve a heterogeneous trajectory $\tau$ for $q$
    \State Compute task reward $r_{\mathrm{task}}$
    \State Obtain inspection scores $\{s_t^{\mathrm{ins}}\}_{t=1}^T$ from $\mathcal{J}(q,\tau)$
    \For{$t=1,\ldots,T$}
      \State Compute $r_t$ using task feedback and inspection scores (Appendix~\ref{app:training})
      \State Add $(x_t,o_t,r_t)$ to $\mathcal{D}_{i_t}$
    \EndFor
  \EndFor
  \For{$i=1,\ldots,N$}
    \State Update $M_i$ with $\mathcal{D}_i$ using reward-weighted optimization
  \EndFor
\EndFor
\State \Return Trained heterogeneous agents $\mathcal{M}=\{M_1,\ldots,M_N\}$
\end{algorithmic}
\end{algorithm}
\endgroup

\section{Interaction and Inspection Protocol}
\label{app:protocol}

This section details response requirements, inspector inputs, and task-specific scoring criteria.

\FloatBarrier
\subsection{Response Schema and Field Requirements}
\label{app:output_schema}
Following the interaction protocol defined in Section~\ref{sec:collaborative_reasoning}, each agent produces a serialized response with the following fields and action-specific requirements at every collaboration step.

\noindent\textbf{Field requirements.}
\begin{itemize}
    \item \texttt{reasoning}: the current step's explanation or verification.
    \item \texttt{tentative\_answer}: the current candidate---a MuSiQue answer, a GSM-Hard number, a MATH expression, a Conifer response, or code in its designated field.
    \item \texttt{action}: the serialized control decision corresponding to $a_t$.
    \item \texttt{handoff\_target}: a permitted recipient ID within the model pool for the next handoff.
    \item \texttt{handoff\_note}: the specific issue the recipient should check.
    \item \texttt{confirmed\_answer}: the final answer for \texttt{terminate}, and null otherwise. For termination, the handoff fields remain null because no recipient is selected.
\end{itemize}
\FloatBarrier
\subsection{Inspector Inputs}

\paragraph{Inspector model.}
Qwen3-14B serves as the shared inspector for all five datasets.

\paragraph{Input visibility.}
Table~\ref{tab:app-inspector-inputs} summarizes inspector inputs across tasks.

\begin{table}[!htbp]
\centering
\caption{Information Available to the Inspector.}
\label{tab:app-inspector-inputs}
\small
\setlength{\tabcolsep}{3pt}
\renewcommand{\arraystretch}{1.15}
\begin{tabular}{@{\hspace{\tabcolsep}}p{0.14\linewidth}p{0.15\linewidth}p{0.15\linewidth}p{0.15\linewidth}p{0.16\linewidth}p{0.15\linewidth}@{\hspace{\tabcolsep}}}
\toprule
\rowcolor{appendixHeader}
\appendixHeading{Information} & \appendixHeading{MuSiQue} & \appendixHeading{GSM-Hard} & \appendixHeading{MATH} & \appendixHeading{MultiPL-E} & \appendixHeading{Conifer} \\
\midrule
Task context & Question and reference paragraphs & Problem statement & Problem statement & Source prefix and language & Instruction and parsed constraints \\
\midrule
\rowcolor{appendixStripe}
Reference information & Gold answer & Gold numeric answer & Gold mathematical answer & No reference code supplied & Reference response for audit \\
\midrule
Candidate and final outputs & Tentative and final answers & Tentative and confirmed answers & Tentative and confirmed answers & Tentative and confirmed code completions & Response drafts and final response \\
\midrule
\rowcolor{appendixStripe}
Agent identity & Agent ID per step & Agent ID per step & Agent ID per step & Agent ID per step & Agent ID per step \\
\midrule
Step sequence & All recorded steps and actions & All recorded steps and actions & All recorded steps and actions & All recorded steps and actions & All recorded steps and actions \\
\midrule
\rowcolor{appendixStripe}
Reasoning text & Up to 400 characters per step & Untruncated reasoning & Untruncated reasoning & Up to 1,200 characters per step & Up to 1,200 characters per step \\
\bottomrule
\end{tabular}
\end{table}

\FloatBarrier
\subsection{Task-Specific Inspection}

The inspector evaluates reasoning quality and the usefulness of each step's interaction decision. Table~\ref{tab:app-rubric} summarizes the scoring criteria for all five datasets used in our experiments.

\begin{table}[!htbp]
\centering
\caption{Scoring Dimensions and Criteria by Dataset.}
\label{tab:app-rubric}
\small
\setlength{\tabcolsep}{3pt}
\renewcommand{\arraystretch}{1.15}
\begin{tabular}{@{\hspace{\tabcolsep}}p{0.14\linewidth}p{0.20\linewidth}p{0.16\linewidth}p{0.42\linewidth}@{\hspace{\tabcolsep}}}
\toprule
\rowcolor{appendixHeader}
\appendixHeading{Dataset} & \appendixHeading{Dimensions} & \appendixHeading{Score range} & \appendixHeading{Criteria} \\
\midrule
MuSiQue & Reasoning; action & $[-1,1]$, continuous & Evidence citations and multi-hop reasoning; supported correction and verification; penalties for regression and redundant interaction. \\
\midrule
\rowcolor{appendixStripe}
GSM-Hard & Process & $[-1,1]$, step size $0.5$ & Correct, independently checkable arithmetic and appropriate actions; penalties for weak verification, misleading computation, and unsafe stopping. Incorrect candidates cannot receive positive scores. \\
\midrule
MATH & Process & $[-1,1]$, step size $0.5$ & Correct mathematical reasoning and useful verification or correction; penalties for misleading work, regression, and unsafe stopping. Incorrect candidates cannot receive positive scores. \\
\midrule
\rowcolor{appendixStripe}
MultiPL-E & Reasoning; action & $[-1,1]$, continuous & Implementation logic, edge cases, syntax, and completion boundaries; useful bug fixes and code audits; penalties for invalid completions and unsupported confirmation. \\
\midrule
Conifer & Reasoning; action; content; final quality & Step: $[-1,1]$; final: $[0,1]$ & Constraint analysis, draft quality, and useful handoff or stopping decisions; penalties for repetition, premature stopping, and regression. Final quality evaluates the submitted response. \\
\bottomrule
\end{tabular}
\end{table}

\FloatBarrier
\section{Training and Inference Details}
\label{app:training}

\subsection{Task Rewards and Training Weights}

Let $S(q,\tau)\in[0,1]$ denote the task score and $s_t^{\mathrm{ins}}\in[-1,1]$ the inspection score. For MuSiQue and Conifer, the training reward is
\begin{equation}
 r_{\mathrm{task}}=2S(q,\tau)-1,
 \qquad
 r_t=\alpha r_{\mathrm{task}}+(1-\alpha)s_t^{\mathrm{ins}}
 \label{eq:app-mixed-reward}
\end{equation}
with $\alpha=0.6$ for MuSiQue and $\alpha=0.5$ for Conifer. The corresponding task-feedback coefficients in the base reward are $\alpha=0.35$ for GSM-Hard, $\alpha=0.35$ for MATH, and $\alpha=0.3$ for MultiPL-E.

\subsection{Agent-Wise Optimization}

For context $x$ and serialized response $o$, let $\ell_{\mathrm{all}}(x,o)$ denote the mean log probability over response tokens. The signed reward-weighted objective follows a score-function policy-gradient form~\citep{williams1992simple}:
\begin{equation}
 \mathcal{L}_i
 =-\mathbb{E}_{(x,o,r)\sim\mathcal{D}_i}\!\left[
 \mathbb{I}[r>0]r\ell_{\mathrm{all}}(x,o)
 +\mathbb{I}[r<0]r\ell_{\mathrm{neg}}(x,o)
 \right]
 +\beta\,\mathbb{E}_{(x,o,r)\sim\mathcal{D}_i}
 \left[\overline{D}_{\mathrm{KL}}(x,o)\right]
 \label{eq:app-rwr}
\end{equation}
where $\overline{D}_{\mathrm{KL}}$ is the mean token-level KL divergence to the frozen SFT reference~\citep{ziegler2019finetuning}.

\subsection{Hardware and LoRA Configuration}

Training uses $8\times$ NVIDIA A800-SXM4-80GB GPUs with a separate LoRA adapter~\citep{hu2022lora} per model and frozen base parameters. The adapters use rank 64, scaling alpha 64, and dropout 0.05. LoRA targets the attention projections \texttt{q\_proj}, \texttt{k\_proj}, \texttt{v\_proj}, and \texttt{o\_proj}, and the MLP projections \texttt{gate\_proj}, \texttt{up\_proj}, and \texttt{down\_proj} in every agent-specific adapter.

\FloatBarrier
\section{Baseline Implementations and Comparison Protocol}
\label{app:baselines}

This section details the implementation and evaluation settings of each baseline. Single-agent controls are Qwen3-14B~\citep{yang2025qwen3} and Qwen3-14B SE-RL. Homogeneous controls are an untrained MAS of three Qwen3-8B agents (Homo. MAS) and a collaboratively trained system of three Qwen3-4B agents (Homo. CL). Established multi-agent comparisons are MAD, AgentVerse, AFlow, and GPTSwarm~\citep{du2024debate,chen2024agentverse,zhang2025aflow,zhuge2024gptswarm}; multi-agent reinforcement-learning comparisons are MAGRPO, MAPoRL, and AT-GRPO~\citep{liu2026collaboration,park2025maporl,zhao2026strongermas}. All use Qwen3 models with the assignments and task-specific configurations detailed below.

\subsection{Single-Agent and Homogeneous Controls}

\paragraph{Qwen3-14B.}
We use the original Qwen3-14B checkpoint as a single-agent baseline for every task.

\paragraph{Qwen3-14B SE-RL.}
We initialize the policy and evaluator from Qwen3-14B, freeze the evaluator during training, and use it to score the policy's generated trajectories.

\paragraph{Homogeneous MAS.}
We use three original Qwen3-8B checkpoints as agents for multi-step collaborative reasoning under the shared interaction protocol used by our system.

\paragraph{Homogeneous collaborative learning (Homo. CL).}
We use the same training and evaluation configuration as CL, replacing its heterogeneous Qwen3-1.7B, Qwen3-4B, and Qwen3-8B model pool with three independently adapted Qwen3-4B agents (12B total).

\subsection{MAD}

MAD organizes reasoning as a multi-round debate in which multiple agents first solve a task independently and then revise their answers after observing the other agents' responses from the preceding round~\citep{du2024debate}. After the final round, the agents' answers are aggregated by majority vote. Our implementation follows MAD's core workflow of independent initial reasoning, iterative peer-conditioned revision, and majority-vote aggregation at test time.

We deploy three debaters using Qwen3-1.7B, Qwen3-4B, and Qwen3-8B, matching the model pool of our heterogeneous MAS. With three total rounds, the nominal cost is nine agent calls per problem, excluding retries. This preserves MAD's three-agent debate and majority-vote structure while making minor adaptations for our Qwen3-based comparison setting.

\subsection{AgentVerse}

AgentVerse recruits three task-specific specialists to generate responses and uses an evaluator to determine whether another iteration is needed~\citep{chen2024agentverse}. The default maximum is three iterations. The evaluator scores the team's responses on a 0--10 scale, with early stopping at a score of at least 8. It receives the task input and agent responses, not a gold answer. The collective answer uses normalized voting with tie handling. This implementation follows AgentVerse's core workflow of expert recruitment, collaborative decision-making, and evaluation-guided iteration.

The three specialist agents use Qwen3-1.7B, Qwen3-4B, and Qwen3-8B, while both the recruiter and evaluator use Qwen3-8B by default. With one recruitment call and four calls per iteration, the nominal range is 5--13 calls, excluding retries. This preserves AgentVerse's separation of specialist, recruitment, and evaluation calls while adapting the backbones and team size to the shared setting.

\subsection{AFlow}

AFlow automates agentic workflow generation by using an optimizer model to propose executable workflows composed of reusable operators, evaluating candidate workflows on task examples, and iteratively refining them through search~\citep{zhang2025aflow}. The selected workflow is then frozen and executed independently at evaluation time. Our implementation follows AFlow's core paradigm of workflow search, execution-based evaluation, and deployment of the selected workflow.

We use Qwen3-14B as the workflow optimizer. During workflow execution, solving calls rotate across Qwen3-1.7B, Qwen3-4B, and Qwen3-8B, while candidate selection and answer extraction use Qwen3-8B by default. Workflow search uses 20 examples and 20 iterations. The optimizer is not invoked during evaluation, preserving AFlow's search-then-execute design.

\subsection{GPTSwarm}

GPTSwarm represents a multi-agent system as an optimizable computational graph, where typed language-model nodes perform specialized operations and directed edges determine how intermediate outputs are exchanged~\citep{zhuge2024gptswarm}. The original method optimizes graph connectivity, while graph execution propagates information toward a designated output node. Our implementation retains GPTSwarm's core graph-structured execution and typed-node composition.

The GPTSwarm-Fixed adaptation executes a fixed seven-node directed acyclic graph: two input/output nodes, three chain-of-thought nodes, one debate node, and one final aggregator. The default routing maps input/output nodes to the 1.7B model, chain-of-thought nodes to the 4B model, and debate/aggregation nodes to the 8B model. Using the original task input, the aggregator synthesizes one answer from its predecessors rather than taking a majority vote.

\subsection{MAGRPO}

We adapt MAGRPO~\citep{liu2026collaboration} to three agents backed by Qwen3-1.7B, Qwen3-4B, and Qwen3-8B, each with an independent LoRA adapter. Training samples groups of eight collaborative episodes and computes episode-relative advantages, with dynamic sampling used to maintain usable groups. The clipped objective uses clip ratio 0.2, advantage clipping at 4, and sequence-mean reduction. The KL coefficient is 0.02 except on MultiPL-E, where it is zero. All three policies use learning rate $5{\times}10^{-6}$. Evaluation runs two synchronous rounds.

\subsection{MAPoRL}

We implement MAPoRL~\citep{park2025maporl} with Qwen3-1.7B, Qwen3-4B, and Qwen3-8B agents and independent LoRA adapters. Each collaborative episode contains three rounds with one proposal from each agent per round. Training uses reward discount 0.3, PPO clip ratio 0.2, KL coefficient 0.002, one PPO epoch, and eight minibatches, with learning rate $10^{-5}$.

\subsection{AT-GRPO}

We use the agent-and-turn grouped optimization procedure of AT-GRPO~\citep{zhao2026strongermas}. The three Qwen3-1.7B, Qwen3-4B, and Qwen3-8B policies have independent LoRA adapters and use spine-tree sampling with three training rounds and group size eight. The clipped objective uses ratio 0.2, KL coefficient 0.02, advantage clipping at 4, and sequence-mean reduction. The A1/A2/A3 learning rates are $2{\times}10^{-5}$, $1.5{\times}10^{-5}$, and $10^{-5}$, respectively. Evaluation uses four synchronous rounds on MuSiQue, GSM-Hard, MATH, and Conifer, and three rounds on MultiPL-E.

\FloatBarrier
\s.}
For each dataset, our method and all baselines are evaluated using the same task-specific evaluatoretcounter{subsection}{9}
\subsection{Conifer Evaluation Metrics}
\label{app:conifer_metrics}

For Conifer, we report two deterministic constraint-compliance metrics computed from each final response. The evaluator first applies a conservative parser to the instruction. For Coverage, it retains numbered requirements containing 2--12 words and at most 120 characters. For Explicit, it extracts mechanically verifiable constraints comprising response presence; requested output formats (bulleted, ordered, or general lists, tables, code, and paragraphs); word, sentence, and list-item limits; and quoted terms that must appear in the response. Neither uses inspectors or references.

Let $a_q$ be the final response for example $q$, $\mathcal{R}_q$ the set of parsed numbered requirements, and $\mathcal{T}(s)$ the set of unique lower-cased, non-stopword alphanumeric tokens in text $s$, allowing apostrophes, underscores, and hyphens within tokens. For an individual requirement $r$, its lexical coverage is
\begin{equation}
 c(a_q,r)=
 \begin{cases}
 1, & |\mathcal{T}(r)|=0,\\
 \dfrac{|\mathcal{T}(r)\cap\mathcal{T}(a_q)|}{|\mathcal{T}(r)|}, & \text{otherwise}.
 \end{cases}
\end{equation}
The per-example Coverage score is
\begin{equation}
 \operatorname{Cov}_q=
 \begin{cases}
 1, & |\mathcal{R}_q|=0,\\
 \dfrac{1}{|\mathcal{R}_q|}\displaystyle\sum_{r\in\mathcal{R}_q}c(a_q,r), & \text{otherwise}.
 \end{cases}
\end{equation}
Thus, Coverage measures the average fraction of each parsed requirement's content tokens that appears in the final response; it is a lexical recall measure rather than a semantic correctness judgment.

For Explicit, let $\mathcal{B}_q$ be the set containing the response-presence check and every mechanically testable constraint extracted for example $q$, and let $b(a_q)\in\{0,1\}$ indicate whether check $b$ passes. The per-example score is
\begin{equation}
 \operatorname{Expl}_q=
 \frac{1}{|\mathcal{B}_q|}\sum_{b\in\mathcal{B}_q} b(a_q).
\end{equation}
All checks receive equal weight. Format and required-term checks use pattern or case-insensitive substring matching, while count constraints compare the measured numbers of words, sentences, or list items with the requested limits. Finally, for $N$ evaluated responses, the reported table entries are $100N^{-1}\sum_q\operatorname{Cov}_q$ and $100N^{-1}\sum_q\operatorname{Expl}_q$. All evaluation items, including unsuccessful predictions, remain in the denominator for every reported aggregate.

\FloatBarrier
\section{Additional Analyses}
\label{app:analysis}

\subsection{Efficiency and Behavioral Metrics}
\label{app:efficiency_metrics}

\paragraph{Amortized parameter-weighted output-token cost.}
We quantify computation cost using amortized parameter-weighted output tokens. This measure combines output tokens generated while solving evaluation problems with an amortized share of the output tokens generated during workflow search. The shared normalization constant is $1.7+4+8+14=27.7$ (in billions), corresponding to the four nominal model sizes used across methods. For a given method and benchmark, the per-problem cost is
\begin{equation}
 C_{\mathrm{out}}(q)
 =\sum_{c\in\mathcal{C}_{\mathrm{eval}}(q)}
   \frac{B_{m(c)}}{27.7}\,L_c^{\mathrm{out}}
 +\frac{1}{N_{\mathrm{eval}}}
  \sum_{c\in\mathcal{C}_{\mathrm{search}}}
   \frac{B_{m(c)}}{27.7}\,L_c^{\mathrm{out}}
 \label{eq:app-weighted-tokens}
\end{equation}
where $B_{m(c)}$ is the nominal parameter count, in billions, of the model producing generation $c$, and $L_c^{\mathrm{out}}$ is the corresponding recorded output-token count. The set $\mathcal{C}_{\mathrm{eval}}(q)$ contains all model generations included during evaluation for problem $q$, including those arising from retries, meta-model evaluations, aggregation, and separate bootstrap solvers. For AFlow, $\mathcal{C}_{\mathrm{search}}$ contains the generations from one recorded workflow search, including those produced by the 14B optimizer and the 1.7B/4B/8B executors used to evaluate candidate workflows. This search cost is divided across the $N_{\mathrm{eval}}$ problems in the evaluation run, not the search examples. Methods without workflow search have $\mathcal{C}_{\mathrm{search}}=\emptyset$. Reported means combine evaluation and search token totals, apply the model weights, and divide by $N_{\mathrm{eval}}$ within each benchmark.

AFlow's 14B optimizer contributes only through the search term, not evaluation-time workflow execution. The metric includes amortized search but excludes training rollouts, teacher generation, training-time inspection, parameter optimization, and offline evaluator replay. It jointly accounts for model size and generated output length. The common normalization does not affect cost ratios, and the amortized search cost per problem decreases as the evaluation set grows.

\paragraph{Reasoning overlap.}
The coordination analysis measures lexical overlap on adjacent steps in which the source agent hands off to the agent that acts next. For receiving agent $i$, the overlap rate is
\begin{equation}
 O_i(\delta)=
 \frac{\sum_{(u,v)\in\mathcal{H}_i}
  \mathbb{I}[\operatorname{ROUGE\!\text{-}\!L}_{F_1}(y_u,y_v)\geq\delta]}
 {|\mathcal{H}_i|}
 \label{eq:app-overlap}
\end{equation}
where $\mathcal{H}_i$ is the set of eligible handoffs received by $i$. This is a threshold-exceedance rate, not a mean similarity score or an embedding-based redundancy metric. Its value depends on the threshold $\delta$, text preprocessing, and the set of valid handoff pairs. Low lexical overlap alone does not prove useful complementary reasoning: an unrelated response can also have low overlap.

\paragraph{Correct-answer discovery and retention.}
Let $\mathcal{V}_q$ be the method-specific set of extracted intermediate candidate answers, and let $g_q(v)$ indicate correctness under the benchmark's grader. Define
\begin{equation}
 D_q=\mathbb{I}[\exists v\in\mathcal{V}_q:g_q(v)=1],
 \qquad
 \mathrm{Discovery}=\frac{\sum_q D_q}{|\mathcal{Q}_{\mathrm{eval}}|},
 \qquad
 \mathrm{Retention}=\frac{\sum_q D_q g_q(\widehat{a}_q)}{\sum_q D_q}
 \label{eq:app-retention}
\end{equation}
Retention is undefined if no correct candidate is discovered. The candidate set $\mathcal{V}_q$ contains intermediate proposals, which are distinct from final aggregation outputs. Differences in extraction coverage can affect cross-method comparisons. A final answer can sometimes be correct even when the extractor finds no correct intermediate candidate, so discovery multiplied by retention need not equal final accuracy. Retention measures preservation of available correct answers; it does not by itself establish the causal mechanism responsible for a performance gain.

\subsection{Self-Evaluation Bias of the 14B Inspector}
\label{app:inspector_bias}

We compare Qwen3-14B inspection scores for heterogeneous-MAS and SE-RL trajectories. Table~\ref{tab:app-inspector-bias} reports mean scores and the share of paired questions favoring the self-generated trajectory.

\begin{table}[!htbp]
\centering
\caption{Qwen3-14B inspection scores for heterogeneous MAS and self-generated SE-RL trajectories. The final column reports the proportion of paired questions on which the self-generated trajectory scores higher than its paired heterogeneous-MAS trajectory.}
\label{tab:app-inspector-bias}
\small
\setlength{\tabcolsep}{4pt}
\renewcommand{\arraystretch}{1.15}
\begin{tabular}{@{\hspace{\tabcolsep}}p{0.20\linewidth}p{0.22\linewidth}p{0.24\linewidth}p{0.24\linewidth}@{\hspace{\tabcolsep}}}
\toprule
\rowcolor{appendixHeader}
\appendixHeading{Dataset} & \appendixHeading{MAS score} & \appendixHeading{Self-generated score} & \appendixHeading{Questions with higher self score} \\
\midrule
MuSiQue & 0.0194 & 0.8692 & 97.25\% \\
\midrule
\rowcolor{appendixStripe}
GSM-Hard & 0.0921 & 0.7991 & 93.26\% \\
\midrule
MultiPL-E-8Lang & 0.6381 & 0.8283 & 90.05\% \\
\midrule
\rowcolor{appendixStripe}
MATH & 0.1054 & 0.4911 & 80.37\% \\
\midrule
Conifer & 0.7643 & 0.8715 & 77.27\% \\
\bottomrule
\end{tabular}
\end{table}

Across all five datasets, the inspector assigns higher mean scores to self-generated trajectories. The paired comparison shows the same asymmetry: self-generated trajectories score higher on 77.27\% to 97.25\% of questions, depending on the dataset in the paired evaluation.

\subsection{Qualitative Analysis of Collaboration Patterns}
\suppressfloats[t]

The two most frequent interaction paths after training are the following:
\begin{equation}
 A_1\rightarrow A_2\rightarrow A_3,
 \qquad
 A_1\rightarrow A_2\rightarrow A_3\rightarrow A_1
\end{equation}
with frequencies of 60.12\% and 34.22\%, respectively. Together they account for 94.34\% of the analyzed trajectories. Both paths involve all three models, while the second returns control to the initial agent after the other two have contributed. This concentration indicates a small set of recurrent interaction patterns. Routing alone does not establish functional specialization, which depends on step content and answer transitions rather than agent order in isolation.

Answer transitions provide a complementary view of these patterns. A wrong-to-correct transition indicates correction, while preserving a correct candidate through later steps indicates retention. A correct-to-wrong transition indicates regression. These transitions do not alone establish independent verification, which also depends on the reasoning content. Scheduler-forced handoffs and minimum-participation rules further constrain the paths independently of model-generated choices.

Section~\ref{app:case_studies} examines these distinctions through matched trajectory comparisons.

\FloatBarrier
\section{Case Studies}
\label{app:case_studies}
\raggedbottom

\colorlet{caseStudyBackground}{appendixHeader!35!white}
\definecolor{caseStudyAccent}{HTML}{887B9C}
\colorlet{caseTranscriptHeader}{appendixHeader}
\colorlet{caseTranscriptMeta}{appendixStripe}
\colorlet{caseTranscriptBody}{white}
\colorlet{caseTranscriptInk}{appendixInk}
\colorlet{caseTranscriptRule}{appendixRule}
\definecolor{caseTranscriptEmphasis}{HTML}{8C5142}
\newcommand{\caseTranscriptHeading}[1]{%
  \rowcolor{caseTranscriptHeader}%
  \textcolor{caseTranscriptInk}{\textbf{Step / action}} &
  \textcolor{caseTranscriptInk}{\textbf{#1}} \\
}
\newsavebox{\caseStudyBoxContent}
\newcommand{\caseStudyBox}[1]{%
  \par\addvspace{6pt}%
  \begingroup
  \sbox{\caseStudyBoxContent}{\fcolorbox{caseStudyBackground}{caseStudyBackground}{#1}}%
  \noindent\usebox{\caseStudyBoxContent}%
  \llap{\makebox[\wd\caseStudyBoxContent][l]{%
    \textcolor{caseStudyAccent}{%
      \rule[-\dp\caseStudyBoxContent]{2pt}{\dimexpr\ht\caseStudyBoxContent+\dp\caseStudyBoxContent\relax}}}}%
  \par\endgroup
  \addvspace{4pt}%
}

The following cases illustrate distinct collaboration behaviors in recorded trajectories. Unless explicitly marked as excerpts or summaries, displayed reasoning fields are reproduced in full, including errors; only typography and emphasis are added. Any omitted steps are identified explicitly. Agent identities and zero-based step indices follow the logs. Arrows denote handoffs, and \texttt{confirm\_stop} denotes termination. Muted brown bold text highlights repeated or incorrect reasoning; black bold text highlights the checks discussed in the analysis.

\subsection{MuSiQue: Revisiting Evidence}
\label{app:case_musique}

We compare five zero-shot steps with three trained steps on the same postal-service question.

\begingroup
\setlength{\fboxsep}{7pt}
\caseStudyBox{%
\begin{minipage}{\dimexpr\linewidth-2\fboxsep-2\fboxrule\relax}
\small
\textbf{Question.} When did the person who first brought a postal service into Umayyad lands become caliph?\\[4pt]
\textbf{Reference:} 661. \qquad \textbf{ID:} \texttt{2hop\_\_16581\_16525}.
\end{minipage}}
\endgroup

\paragraph{Zero-shot: repeated reasoning across handoffs.}
\textcolor{caseTranscriptEmphasis}{\textbf{Muted brown bold text}} marks repeated spans. The expansion claim recurs in all five steps; steps 1--4 share the same final sentence, and steps 2 and 4 have identical reasoning. These highlights show textual repetition rather than factual correctness.

\begingroup
\small
\setlength{\tabcolsep}{4pt}
\setlength{\LTcapwidth}{\linewidth}
\setlength{\LTleft}{0pt}
\setlength{\LTright}{\fill}
\setlength{\LTpre}{6pt}
\setlength{\LTpost}{4pt}
\renewcommand{\arraystretch}{1.08}
\arrayrulecolor{caseTranscriptRule}
\begin{longtable}{@{\hspace{3.5pt}}
  >{\columncolor{caseTranscriptMeta}[3.5pt][\tabcolsep]\leavevmode\color{caseTranscriptInk}}p{0.16\linewidth}
  >{\columncolor{caseTranscriptBody}[\tabcolsep][3.5pt]}p{0.80\linewidth}@{\hspace{3.5pt}}}
\caption{MuSiQue, zero-shot MAS: all five steps repeat the same incorrect candidate, 711.}\label{tab:case_musique_zero}\\
\toprule
\caseTranscriptHeading{Original reasoning (emphasis added)}
\midrule
\endfirsthead
\toprule
\caseTranscriptHeading{Original reasoning (emphasis added)}
\midrule
\endhead
\endfoot
\bottomrule
\endlastfoot
0 / $A_1$\newline $\to A_2$ & ``The question asks when the person who first brought postal service into Umayyad lands became caliph. From paragraph [12], it states that the \textcolor{caseTranscriptEmphasis}{\textbf{Umayyad Caliphate started expanding rapidly in the peninsula beginning in 711}}, and the west of the empire gained its independence under Abd-ar-Rahman I with the establishment of the Emirate of C\'ordoba. The Emirate of C\'ordoba was established in 711, and Abd-ar-Rahman I was the first to bring postal service into Umayyad lands. Therefore, the person who first brought postal service into Umayyad lands became caliph in 711.'' \\
\midrule
1 / $A_2$\newline $\to A_3$ & ``The question asks when the person who first brought postal service into Umayyad lands became caliph. Paragraph [12] states that the \textcolor{caseTranscriptEmphasis}{\textbf{Umayyad Caliphate started expanding rapidly in the peninsula beginning in 711}}, and the west of the empire gained its independence under Abd-ar-Rahman I with the establishment of the Emirate of C\'ordoba. \textcolor{caseTranscriptEmphasis}{\textbf{Abd-ar-Rahman I is noted as the first to bring postal service into Umayyad lands, and the Emirate of C\'ordoba was established in 711.}}'' \\
\midrule
2 / $A_3$\newline $\to A_1$ & ``Paragraph [12] states that the \textcolor{caseTranscriptEmphasis}{\textbf{Umayyad Caliphate started expanding rapidly in the peninsula beginning in 711}}, and the west of the empire gained its independence under Abd-ar-Rahman I with the establishment of the Emirate of C\'ordoba. \textcolor{caseTranscriptEmphasis}{\textbf{Abd-ar-Rahman I is noted as the first to bring postal service into Umayyad lands, and the Emirate of C\'ordoba was established in 711.}}'' \\
\midrule
3 / $A_1$\newline $\to A_2$ & ``The question asks when the person who first brought postal service into Umayyad lands became caliph. From paragraph [12], it states that the \textcolor{caseTranscriptEmphasis}{\textbf{Umayyad Caliphate started expanding rapidly in the peninsula beginning in 711}}, and the west of the empire gained its independence under Abd-ar-Rahman I with the establishment of the Emirate of C\'ordoba. \textcolor{caseTranscriptEmphasis}{\textbf{Abd-ar-Rahman I is noted as the first to bring postal service into Umayyad lands, and the Emirate of C\'ordoba was established in 711.}}'' \\
\midrule
4 / $A_2$\newline \texttt{confirm\_stop} & ``Paragraph [12] states that the \textcolor{caseTranscriptEmphasis}{\textbf{Umayyad Caliphate started expanding rapidly in the peninsula beginning in 711}}, and the west of the empire gained its independence under Abd-ar-Rahman I with the establishment of the Emirate of C\'ordoba. \textcolor{caseTranscriptEmphasis}{\textbf{Abd-ar-Rahman I is noted as the first to bring postal service into Umayyad lands, and the Emirate of C\'ordoba was established in 711.}}'' \\
\end{longtable}
\arrayrulecolor{black}
\noindent\textbf{Outcome:} 5 steps; 4 handoffs; final answer 711; EM = 0, F1 = 0.
\par\endgroup

\noindent\small\textbf{Identical handoff note at steps 0--3:} ``Please verify that the person who first brought postal service into Umayyad lands became caliph in 711 using paragraph 12.''\par\normalsize

\paragraph{After training: revisiting the requested relation.}
The trained trajectory initially proposes two incorrect candidates. At the final step, $A_3$ distinguishes \emph{introducing} postal service from \emph{extending} it. Bold highlights mark the question-specific check and conclusion while retaining surrounding errors.

\begingroup
\small
\setlength{\tabcolsep}{4pt}
\setlength{\LTcapwidth}{\linewidth}
\setlength{\LTleft}{0pt}
\setlength{\LTright}{\fill}
\setlength{\LTpre}{6pt}
\setlength{\LTpost}{4pt}
\renewcommand{\arraystretch}{1.08}
\arrayrulecolor{caseTranscriptRule}
\begin{longtable}{@{\hspace{3.5pt}}
  >{\columncolor{caseTranscriptMeta}[3.5pt][\tabcolsep]\leavevmode\color{caseTranscriptInk}}p{0.16\linewidth}
  >{\columncolor{caseTranscriptBody}[\tabcolsep][3.5pt]}p{0.80\linewidth}@{\hspace{3.5pt}}}
\caption{MuSiQue, trained MAS: the final agent distinguishes the introduction of postal service from its later expansion, successfully correcting the earlier answer and arriving at the correct date.}\label{tab:case_musique_trained}\\
\toprule
\caseTranscriptHeading{Original reasoning and candidate (emphasis added)}
\midrule
\endfirsthead
\toprule
\caseTranscriptHeading{Original reasoning and candidate (emphasis added)}
\midrule
\endhead
\endfoot
\bottomrule
\endlastfoot
0 / $A_1$\newline $\to A_2$ & ``Muhammad bin Qasim (paragraph 2) was the first to bring Islam to Punjab, and he was the Umayyad Caliphate's representative in the region. His appointment as a governor and his role in the Umayyad Caliphate's expansion (paragraph 2) implies he was a key figure in the Umayyad administration in the region. Thus, his role in the Umayyad Caliphate's establishment in the region (and not as a caliph in the broader sense) is tied to the Umayyad Caliphate's expansion.''\newline \textit{Candidate: Muhammad bin Qasim.} \\
\midrule
1 / $A_2$\newline $\to A_3$ & ``Paragraph 2 states Muhammad bin Qasim introduced Islam to Punjab under Umayyad rule. Paragraph 8 confirms Umayyad Caliph Abd al-Malik extended postal service. Since Abd al-Malik's reign (723--743) is explicitly mentioned, his postal reforms likely occurred during his caliphate.''\newline \textit{Candidate: 723--743.} \\
\midrule
2 / $A_3$\newline \texttt{confirm\_stop} & ``Paragraph 8 states the Umayyad Caliph Abd al-Malik extended the postal service. Paragraph 20 mentions Muawiyah I's reign (661--80) and his role in establishing the Umayyad dynasty. However, paragraph 8 specifies Abd al-Malik's caliphate (723--743) as the period when the postal service was expanded. \textbf{The question asks for the person who first introduced postal service into Umayyad lands.} Paragraph 8 attributes the introduction of postal service to Muawiyah (paragraph 20's 'Muawiyah I' is likely the same as paragraph 8's 'Muawiyah' who was the first Umayyad caliph, but paragraph 8 says Abd al-Malik 'extended' it. However, paragraph 8 also mentions 'Muawiyah introduced postal service' (paragraph 8: 'Mu'awiyah introduced postal service, Abd al-Malik extended it...'). \textbf{This indicates Muawiyah was the first to introduce it, and his caliphate began in 661 (paragraph 20).}''\newline \textit{Candidate: 661.} \\
\end{longtable}
\arrayrulecolor{black}
\noindent\textbf{Outcome:} 3 steps; 2 handoffs; final answer 661; EM = 1, F1 = 1.
\par\endgroup

\begingroup
\setlength{\fboxsep}{7pt}
\caseStudyBox{%
\begin{minipage}{\dimexpr\linewidth-2\fboxsep-2\fboxrule\relax}
\small
\textbf{Evidence check.} The reference chain links [8], ``Mu'awiyah introduced postal service, Abd al-Malik extended it throughout his empire,'' to [1], ``Caliph Muawiyah (661--80)''. The zero-shot agents instead cite [12], whose supplied excerpt describes expansion into Portugal beginning in 711, not the introduction of postal service. Trained $A_3$ cites [20], which connects Muawiyah I and the beginning of the dynasty to events in 661, rather than the reference paragraph [1], while still supporting the correct accession date.
\end{minipage}}
\endgroup

\paragraph{Repetition versus substantive re-examination.}
Although the trained trajectory begins with an incorrect candidate, subsequent agents do not simply reaffirm the initial hypothesis. Instead, the focus shifts from an irrelevant entity to postal-service evidence before the final agent distinguishes the service's introduction from its later expansion and identifies the correct person--date relation. This progression shows how cross-agent re-examination can recover from an incorrect initial direction, unlike the repetitive confirmation in the zero-shot trajectory.

\paragraph{Successful correction through cross-perspective collaboration.}
The final agent brings a complementary perspective to the shared reasoning by examining whether the cited evidence establishes the specific relation requested by the question. Distinguishing who first introduced postal service from who later extended it allows the agent to identify Muawiyah and return the correct accession date, 661. Here, cross-model collaboration enables a later model to revisit the interpretation of shared evidence rather than simply continue the preceding line of reasoning. The successful final correction illustrates how complementary reasoning perspectives can be translated into a more accurate answer through explicit cross-model interaction in the final response.

\FloatBarrier

\noindent\begin{minipage}{\linewidth}
\subsection{MuSiQue: Questioning Prior Reasoning}
\label{app:case_musique_question}

This second MuSiQue example concerns the interpretation of the question rather than the distinction between two historical relations. We show the zero-shot trajectory's first and final steps and all three trained steps; the omitted zero-shot steps are described below.

\begingroup
\setlength{\fboxsep}{7pt}
\caseStudyBox{%
\begin{minipage}{\dimexpr\linewidth-2\fboxsep-2\fboxrule\relax}
\small
\textbf{Question.} When is Celebrity Big Brother coming to the network which, along with ABC and the network which hosts Westinghouse Playhouse, is the other major broadcaster based in NY?\\[4pt]
\textbf{Reference:} February 7, 2018.\\[2pt]
\textbf{ID:} \texttt{3hop1\_\_140554\_2053\_52946}.
\end{minipage}}
\endgroup
\end{minipage}\par

\paragraph{Zero-shot: retaining an incomplete task interpretation.}
Although the question asks \emph{when} the program will air, the agents repeatedly treat it as asking for a network. They also select NBC rather than excluding it alongside ABC. All four candidates are NBC. The complete reasoning texts of the omitted steps 1 and 2 are identical word for word to that of the final step shown in Table~\ref{tab:case_musique_question_zero}, including punctuation and spacing. Both omitted steps hand off again without resolving the date.

\begingroup
\small
\setlength{\tabcolsep}{4pt}
\setlength{\LTcapwidth}{\linewidth}
\setlength{\LTleft}{0pt}
\setlength{\LTright}{\fill}
\setlength{\LTpre}{6pt}
\setlength{\LTpost}{4pt}
\renewcommand{\arraystretch}{1.08}
\arrayrulecolor{caseTranscriptRule}
\begin{longtable}{@{\hspace{3.5pt}}
  >{\columncolor{caseTranscriptMeta}[3.5pt][\tabcolsep]\leavevmode\color{caseTranscriptInk}}p{0.16\linewidth}
  >{\columncolor{caseTranscriptBody}[\tabcolsep][3.5pt]}p{0.80\linewidth}@{\hspace{3.5pt}}}
\caption{MuSiQue, zero-shot MAS: the first and final steps answer which network rather than when the program will air.}\label{tab:case_musique_question_zero}\\
\toprule
\caseTranscriptHeading{Original reasoning and candidate (emphasis added)}
\midrule
\endfirsthead
\toprule
\caseTranscriptHeading{Original reasoning and candidate (emphasis added)}
\midrule
\endhead
\endfoot
\bottomrule
\endlastfoot
0 / $A_1$\newline $\to A_2$ & ``\textcolor{caseTranscriptEmphasis}{\textbf{The question asks for the network}} that is the other major broadcaster based in New York, along with ABC and the network that hosts Westinghouse Playhouse. From the paragraphs, we know that ABC is one of the major broadcasters in New York. The network that hosts Westinghouse Playhouse is NBC. Therefore, the other major broadcaster based in NY is NBC.''\newline \textit{Candidate: NBC.} \\
\midrule
3 / $A_1$\newline \texttt{confirm\_stop} & ``\textcolor{caseTranscriptEmphasis}{\textbf{The question asks for the network}} that is the other major broadcaster based in New York, along with ABC and the network that hosts Westinghouse Playhouse. From paragraph 1, ABC is a major broadcaster in New York. From paragraph 16, Westinghouse Playhouse aired on NBC. Therefore, the other major broadcaster based in NY is NBC.''\newline \textit{Candidate: NBC.} \\
\end{longtable}
\arrayrulecolor{black}
\noindent\textbf{Outcome:} 4 steps; 3 handoffs; final answer NBC; EM = 0, F1 = 0.
\par\endgroup

\paragraph{After training: returning to the original question.}
The trained candidates progress from CBS to the 2017--2018 season and finally to February 7, 2018. The first two omit the requested premiere date; bold text in Table~\ref{tab:case_musique_question_trained} marks the final agent's return to the temporal question.

\begingroup
\small
\setlength{\tabcolsep}{4pt}
\setlength{\LTcapwidth}{\linewidth}
\setlength{\LTleft}{0pt}
\setlength{\LTright}{\fill}
\setlength{\LTpre}{6pt}
\setlength{\LTpost}{4pt}
\renewcommand{\arraystretch}{1.08}
\arrayrulecolor{caseTranscriptRule}
\begin{longtable}{@{\hspace{3.5pt}}
  >{\columncolor{caseTranscriptMeta}[3.5pt][\tabcolsep]\leavevmode\color{caseTranscriptInk}}p{0.16\linewidth}
  >{\columncolor{caseTranscriptBody}[\tabcolsep][3.5pt]}p{0.80\linewidth}@{\hspace{3.5pt}}}
\caption{MuSiQue, trained MAS: three steps with candidates and handoff notes, ending in a revision of the task interpretation.}\label{tab:case_musique_question_trained}\\
\toprule
\caseTranscriptHeading{Original reasoning and candidate (emphasis added)}
\midrule
\endfirsthead
\toprule
\caseTranscriptHeading{Original reasoning and candidate (emphasis added)}
\midrule
\endhead
\endfoot
\bottomrule
\endlastfoot
0 / $A_1$\newline $\to A_2$ & ``Celebrity Big Brother is associated with CBS (paragraph 19). Westinghouse Playhouse is on NBC (paragraph 16). The question asks for the network that is not ABC and is not the same as CBS, but is a major network in New York. The correct answer is CBS, which is not explicitly excluded from the 'major broadcasters in New York' list, but the question implies a network that is not ABC and not the same as NBC (since Westinghouse Playhouse is on NBC). Thus, the network in question is CBS.''\newline \textit{Candidate: CBS.}\newline \textbf{Handoff note:} ``Verify that CBS is the correct network for the season of Celebrity Big Brother, and that it is not the same as ABC or NBC.'' \\
\midrule
1 / $A_2$\newline $\to A_3$ & ``Paragraph 19 states Celebrity Big Brother aired on CBS during 2017-2018. Paragraph 16 confirms Westinghouse Playhouse (NBC) is a New York network. Paragraph 4 notes ABC and CBS are both NYC-based. Thus, CBS is the network alongside ABC and NBC.''\newline \textit{Candidate: 2017-2018.}\newline \textbf{Handoff note:} ``Check if CBS's 2017-2018 schedule aligns with the NYC-based networks' overlap.'' \\
\midrule
2 / $A_3$\newline \texttt{confirm\_stop} & ``Paragraph 17 mentions Celebrity Big Brother (American TV series) will air on CBS during the 2017-18 network season. Paragraph 4 states ABC and CBS are both headquartered in NYC. Paragraph 16 notes Westinghouse Playhouse aired on NBC. The question asks for the network (other than ABC and NBC) based in NY. However, the phrasing implies the network that hosts Celebrity Big Brother (CBS) is being asked, but the user might have intended to ask about a different network. Wait, the question says 'along with ABC and the network which hosts Westinghouse Playhouse (NBC), is the other major broadcaster based in NY'. The three major NYC networks are ABC, CBS, and NBC (paragraph 4). The question seems to ask when Celebrity Big Brother (CBS) is coming to the network that is not ABC or NBC, but this is impossible. However, the correct interpretation is that the question is asking for the network (CBS) that, along with ABC and NBC (hosting Westinghouse Playhouse), is a major NYC-based network. \textbf{But the actual question is about when Celebrity Big Brother (CBS) is coming to the network (CBS) that is part of the NYC-based trio.} This seems redundant. Alternatively, the question might have a typo, but based on the given data, the answer refers to the CBS season premiere date. Paragraph 19 states the 2017-18 season premiered on February 7, 2018.''\newline \textit{Candidate: February 7, 2018.} \\
\end{longtable}
\arrayrulecolor{black}
\noindent\textbf{Outcome:} 3 steps; 2 handoffs; final answer February 7, 2018; EM = 1, F1 = 1.
\par\endgroup

\begingroup
\setlength{\fboxsep}{7pt}
\caseStudyBox{%
\begin{minipage}{\dimexpr\linewidth-2\fboxsep-2\fboxrule\relax}
\small
\textbf{Evidence check.} Paragraph [16] identifies NBC as the network on which \emph{Westinghouse Playhouse} aired. Paragraph [4] lists ABC, CBS, and NBC as the three major networks headquartered in New York, identifying CBS as the remaining network. Paragraph [17] states that \emph{Celebrity Big Brother} will air on CBS and that ``CBS announced that the series is set to premiere on February 7, 2018 and conclude on February 25, 2018.'' The final model output attributes the date to [19]; this citation is retained as logged, while [17] supplies the date in the provided evidence excerpts used for comparison.
\end{minipage}}
\endgroup

\paragraph{Questioning and revising prior reasoning.}
The distinction here is between identifying an intermediate entity and answering the original question. Whereas the zero-shot agents remain focused on network identification, the final trained agent explicitly recognizes that the question asks when the program will air. It retains CBS as the relevant network but replaces the preceding season-level candidate with the premiere date. This illustrates revision of the task interpretation, not merely substitution of one network name for another. In this case, a later model brings a different perspective to the shared reasoning, correcting the earlier focus on network identification rather than carrying that incomplete interpretation into the reported final answer in the evaluation results.

\FloatBarrier

\noindent\begin{minipage}{\linewidth}
\subsection{GSM-Hard: Verifying a Correct Answer}
\label{app:case_gsm_hard}

This example shows correct-answer retention through an alternative arithmetic check. We show both recorded steps and the handoff note, with mathematical expressions typeset.

\begingroup
\setlength{\fboxsep}{7pt}
\caseStudyBox{%
\begin{minipage}{\dimexpr\linewidth-2\fboxsep-2\fboxrule\relax}
\small
\textbf{Question.} If there are 657881 eggs in a basket, and there are twice as many eggs in a second basket, how many eggs are in both baskets put together?\\[4pt]
\textbf{Reference:} 1973643. \qquad \textbf{ID:} \texttt{gsm\_00016}.
\end{minipage}}
\endgroup
\end{minipage}\par

\begin{table}[!htbp]
\centering
\small
\setlength{\tabcolsep}{4pt}
\renewcommand{\arraystretch}{1.08}
\caption{GSM-Hard: $A_2$ verifies $A_3$'s answer via multiplication and place-value decomposition.}
\label{tab:case_gsm_hard}
\arrayrulecolor{caseTranscriptRule}
\begin{tabular}{@{\hspace{3.5pt}}
  >{\columncolor{caseTranscriptMeta}[3.5pt][\tabcolsep]\leavevmode\color{caseTranscriptInk}}p{0.16\linewidth}
  >{\columncolor{caseTranscriptBody}[\tabcolsep][3.5pt]}p{0.80\linewidth}@{\hspace{3.5pt}}}
\toprule
\caseTranscriptHeading{Original reasoning and candidate (emphasis added)}
\midrule
0 / $A_3$\newline $\to A_2$ & ``First, calculate the eggs in the second basket: $657881 * 2 = 1315762$. Then add both baskets: $657881 + 1315762 = 1973643$.''\newline \textit{Candidate: 1973643.}\newline \textbf{Handoff note:} ``Verify the addition step: $657881 + 1315762$'' \\
\midrule
1 / $A_2$\newline \texttt{confirm\_stop} & ``\textbf{Total eggs = $\boldsymbol{3 \times 657881}$.} Breakdown: $600000\times3=1800000$, $50000\times3=150000$, $7000\times3=21000$, $800\times3=2400$, $80\times3=240$, $1\times3=3$. Sum: $1800000+150000=1950000$; $+21000=1971000$; $+2400=1973400$; $+240=1973640$; $+3=1973643$.''\newline \textit{Candidate: 1973643.} \\
\bottomrule
\end{tabular}
\arrayrulecolor{black}

\smallskip
\raggedright\small\textbf{Outcome:} 2 steps; 1 handoff; final answer 1973643; EM = 1.
\end{table}

\paragraph{Preserving a correct answer through verification.}
The receiving agent builds on the preceding reasoning by expressing the same total as three times the first basket's contents and verifying the multiplication through a place-value decomposition. This alternative calculation reaches the same answer while providing an additional computational basis for accepting it. The case illustrates that cross-model collaboration is valuable not only for revising answers, but also for substantiating correct ones: a later model can contribute a complementary reasoning perspective that strengthens the basis for retaining an existing solution. Diverse perspectives thus let agents cross-check the same problem through complementary calculations, improving final-answer reliability.

\FloatBarrier

\subsection{GSM-Hard: Recovering from a Regression}
\label{app:case_gsm_regression}

This second GSM-Hard example shows a correct answer becoming incorrect before being restored at termination. All three steps and the two handoff requests are shown. The highlighted spans contrast the intermediate quantity assumption with the final conversion from legs to pairs.

\begingroup
\setlength{\fboxsep}{7pt}
\caseStudyBox{%
\begin{minipage}{\dimexpr\linewidth-2\fboxsep-2\fboxrule\relax}
\small
\textbf{Question.} Mario needs to buy snowshoes for his 8374485 sled dogs. Assuming his dogs each has four legs and each pair of snowshoes costs \$12.00, how much will it cost him to buy snowshoes for all of his dogs?\\[4pt]
\textbf{Reference:} 200987640. \qquad \textbf{ID:} \texttt{gsm\_00076}.
\end{minipage}}
\endgroup
\par

\begingroup
\small
\setlength{\tabcolsep}{4pt}
\setlength{\LTcapwidth}{\linewidth}
\setlength{\LTleft}{0pt}
\setlength{\LTright}{\fill}
\setlength{\LTpre}{6pt}
\setlength{\LTpost}{4pt}
\renewcommand{\arraystretch}{1.08}
\arrayrulecolor{caseTranscriptRule}
\begin{longtable}{@{\hspace{3.5pt}}
  >{\columncolor{caseTranscriptMeta}[3.5pt][\tabcolsep]\leavevmode\color{caseTranscriptInk}}p{0.16\linewidth}
  >{\columncolor{caseTranscriptBody}[\tabcolsep][3.5pt]}p{0.80\linewidth}@{\hspace{3.5pt}}}
\caption{GSM-Hard: $A_1$ corrects $A_2$'s one-pair-per-dog regression, restoring the initial answer.}\label{tab:case_gsm_regression}\\
\toprule
\caseTranscriptHeading{Original reasoning and answer (emphasis added)}
\midrule
\endfirsthead
\toprule
\caseTranscriptHeading{Original reasoning and answer (emphasis added)}
\midrule
\endhead
\endfoot
\bottomrule
\endlastfoot
0 / $A_3$ (8B)\newline $\to A_2$ & ``First, calculate the total number of snowshoe pairs needed: 8,374,485 dogs * 2 pairs per dog = 16,748,970 pairs. Then multiply by the cost per pair: 16,748,970 * \$12.00 = \$200,987,640.''\newline \textit{Candidate: 200987640.}\newline \textbf{Handoff note:} ``Verify the multiplication step: 16,748,970 * 12'' \\
\midrule
1 / $A_2$ (4B)\newline $\to A_1$ & ``First, calculate the total number of snowshoe pairs needed: 8,374,485 dogs * \textcolor{caseTranscriptEmphasis}{\textbf{1 pair per dog}} = 8,374,485 pairs. Then multiply by the cost per pair: 8,374,485 * \$12.00 = \$100,493,820.''\newline \textit{Candidate: 100493820.}\newline \textbf{Handoff note:} ``Verify the multiplication step: 8,374,485 * 12'' \\
\midrule
2 / $A_1$ (1.7B)\newline \texttt{confirm\_stop} & ``First, determine the number of snowshoe pairs per dog: \textbf{4 legs / 2 legs per pair = 2 pairs.} Total pairs = 8,374,485 dogs * 2 pairs/dog = 16,748,970 pairs. Total cost = 16,748,970 pairs * \$12.00/pair = \$200,987,640.''\newline \textit{Confirmed answer: 200987640.} \\
\end{longtable}
\arrayrulecolor{black}
\noindent\textbf{Outcome:} 3 steps; 2 handoffs; final answer 200987640; EM = 1.
\par\endgroup

\noindent\small\textbf{Note:} ``Independently reverified the original problem's arithmetic and confirmed the final numeric answer.''\par\normalsize

\paragraph{Recovering from an intermediate reasoning regression.}
The initial agent correctly assigns two pairs of snowshoes to each dog. The intermediate agent changes this quantity to one pair, halving the total cost. Its multiplication is correct for that quantity; the error lies in the number of pairs per dog, not in the multiplication itself. The final agent returns to the problem's four-leg condition and explicitly derives two pairs per dog before recomputing the total. This check repairs the incorrect premise introduced in the preceding step and restores the original correct answer. Unlike the preceding GSM-Hard case, where correctness is retained throughout, this trajectory demonstrates recovery after an intervening regression. Here, the repair is made by the smallest agent, $A_1$ (1.7B). Although introducing another model's perspective may yield a reasoning step that does not fully account for the problem's conditions, a subsequent model can revisit that step from a different perspective and bring the reasoning back into alignment with those conditions. In this case, cross-model collaboration enables the pairs-per-dog assumption to be reassessed and corrected before the final answer, illustrating the value of successive cross-perspective checks.

\FloatBarrier

\subsection{MultiPL-E: Completing a Missing Validation Check}
\label{app:case_multiple_validation}

This C++ example concerns a requirement acknowledged in the initial reasoning but omitted from the implementation. Here, $A_1$, $A_2$, and $A_3$ correspond to Qwen3-8B, Qwen3-4B, and Qwen3-1.7B, respectively. All three steps appear below: steps 0--1 use supplied reasoning excerpts, and the final step is summarized. Ellipses mark omitted text. Only the first handoff note is excerpted; the second note and final-step reasoning are not reproduced in the table.

\begingroup
\setlength{\fboxsep}{7pt}
\caseStudyBox{%
\begin{minipage}{\dimexpr\linewidth-2\fboxsep-2\fboxrule\relax}
\small
\textbf{Task.} Given the lengths of the three sides of a triangle, return its area rounded to two decimal places if the triangle is valid; otherwise return $-1$. A valid triangle requires the sum of any two sides to be greater than the third.\\[4pt]
\textbf{Signature:} \texttt{float triangle\_area(long a, long b, long c)}\\[4pt]
\textbf{ID:} \texttt{humaneval/cpp/HumanEval\_71\_triangle\_area}.\\[2pt]
\textbf{Evaluation:} Nine assertions; invalid triangles $(1,2,10)$, $(1,2,3)$, $(2,6,3)$, and $(2,2,10)$ must return $-1$.
\end{minipage}}
\endgroup
\par

\begingroup
\small
\setlength{\tabcolsep}{4pt}
\setlength{\LTcapwidth}{\linewidth}
\setlength{\LTleft}{0pt}
\setlength{\LTright}{\fill}
\setlength{\LTpre}{6pt}
\setlength{\LTpost}{4pt}
\renewcommand{\arraystretch}{1.08}
\arrayrulecolor{caseTranscriptRule}
\begin{longtable}{@{\hspace{3.5pt}}
  >{\columncolor{caseTranscriptMeta}[3.5pt][\tabcolsep]\leavevmode\color{caseTranscriptInk}}p{0.16\linewidth}
  >{\columncolor{caseTranscriptBody}[\tabcolsep][3.5pt]}p{0.80\linewidth}@{\hspace{3.5pt}}}
\caption{MultiPL-E: $A_2$ adds the validation check missing from $A_1$'s implementation, and $A_3$ retains the repair. Reasoning and handoff text is excerpted; the final step is summarized.}\label{tab:case_multiple_validation}\\
\toprule
\caseTranscriptHeading{Reasoning excerpts and candidate code (emphasis added)}
\midrule
\endfirsthead
\toprule
\caseTranscriptHeading{Reasoning excerpts and candidate code (emphasis added)}
\midrule
\endhead
\endfoot
\bottomrule
\endlastfoot
0 / $A_1$ (8B)\newline $\to A_2$ & ``The function needs to check if the three sides form a valid triangle. If they do, it should calculate the area using Heron's formula\ldots\ Otherwise, return -1.''\newline
\textbf{Candidate body:}\newline
\texttt{float s = (a + b + c) / 2.0f;}\newline
\texttt{float area = sqrt(s * (s - a) * (s - b) * (s - c));}\newline
\texttt{return round(area * 100) / 100;}\newline
\textbf{Handoff excerpt:} ``Verify the calculation\ldots\ \textbf{Check if the validation of the triangle sides is properly handled before the calculation.}'' \\
\midrule
1 / $A_2$ (4B)\newline $\to A_3$ & ``The candidate code correctly implements Heron's formula\ldots\ However, \textbf{the validation of the triangle sides is missing.} The function should first check if the sides form a valid triangle before proceeding.''\newline
\textbf{Candidate body:}\newline
\texttt{\textbf{if (a + b > c \&\& a + c > b \&\& b + c > a) \{}}\newline
\hspace*{1em}\texttt{float s = (a + b + c) / 2.0f;}\newline
\hspace*{1em}\texttt{float area = sqrt(s * (s - a) * (s - b) * (s - c));}\newline
\hspace*{1em}\texttt{return round(area * 100) / 100;}\newline
\texttt{\}}\newline
\texttt{\textbf{return -1L;}} \\
\midrule
2 / $A_3$ (1.7B)\newline \texttt{confirm\_stop} & \textbf{Summary (not a quotation):} The final candidate retains the preceding implementation, compressed onto one line, including the triangle-inequality check and the fallback return value. The agent terminates the trajectory. \\
\end{longtable}
\arrayrulecolor{black}
\noindent\textbf{Outcome:} 3 steps; 2 handoffs; the final candidate passes all nine supplied assertions in post-hoc execution.
\par\endgroup

\paragraph{Post-hoc execution check.}
Each step's candidate was separately compiled and run against the supplied tests after the trajectory was collected. All three candidates compiled successfully, but only the latter two passed all nine assertions reported in Table~\ref{tab:case_multiple_execution}.

\begin{table}[!htbp]
\centering
\small
\setlength{\tabcolsep}{4pt}
\renewcommand{\arraystretch}{1.08}
\caption{Post-hoc results for the three C++ candidates on the supplied nine-assertion test suite.}
\label{tab:case_multiple_execution}
\begin{tabular}{@{}p{0.09\linewidth}p{0.16\linewidth}p{0.15\linewidth}p{0.54\linewidth}@{}}
\toprule
\rowcolor{appendixHeader}
\appendixHeading{Step} & \appendixHeading{Agent} & \appendixHeading{Compilation} & \appendixHeading{Test result} \\
\midrule
0 & $A_1$ (8B) & Success & Fails the $-1$ assertion for $(1,2,10)$ and aborts. \\
\rowcolor{appendixStripe}
1 & $A_2$ (4B) & Success & Passes all nine assertions. \\
2 & $A_3$ (1.7B) & Success & Passes all nine assertions. \\
\bottomrule
\end{tabular}
\end{table}

\begingroup
\setlength{\fboxsep}{7pt}
\caseStudyBox{%
\begin{minipage}{\dimexpr\linewidth-2\fboxsep-2\fboxrule\relax}
\small
\textbf{Repair mechanism.} For $(a,b,c)=(1,2,10)$, the initial implementation computes $s=6.5$ and $s-c=-3.5$, making the argument to \texttt{sqrt} negative. The resulting NaN does not satisfy the required return value of $-1$. The added guard rejects this input before evaluating the area. It also rejects the degenerate triangle $(1,2,3)$, for which the initial implementation returns zero rather than $-1$.
\end{minipage}}
\endgroup

\paragraph{Completing an implementation rather than replacing a solution.}
The initial agent notes but omits triangle validation. The recipient checks the candidate against the specification, identifies the missing condition, and adds a validation branch while preserving the area calculation. This complementary review fills a gap in the initial code, showing that cross-model collaboration can check and complete an existing solution. The final model retains the repair, and the failure-to-pass transition provides executable evidence that it fixes the defect.

\paragraph{A targeted handoff and a corresponding repair.}
The first handoff explicitly requests a check of triangle-side validation, and the next agent addresses precisely that gap. This alignment illustrates how a targeted request can accompany a substantive cross-model contribution. Together, the trajectory and execution results illustrate how complementary reasoning perspectives across models can be translated into a concrete improvement in the final implementation.

\FloatBarrier
